% ------------------------------------------------------------------------
%
\begin{frame}{Binary Search Trees}
  \label{frame:balanced}

  Searching in a binary tree can be costly because, in the worst case,
  the whole tree must be traversed.

  \bigskip

  To improve upon the worst case, two scenarios are desirable:
  \begin{enumerate}

    \item the binary tree should be as \textbf{balanced} as possible,
      that is, the nodes should lie as close as possible to the root
      (for a given notion of distance), and

    \item at a given node, \textbf{only the left or the right subtree
      should be visited}, a decision taken solely by comparing the key
      at the root and the sought key.

  \end{enumerate}

  The simplest solution consists in satisfying the second goal, and
  see later how to also reach the first.

\end{frame}

% ------------------------------------------------------------------------
%
\begin{frame}{Binary Search Trees}

  A \textbf{binary search tree} is a binary tree where the key of the
  roo (if any) is greater than or equal to the nodes in the left
  subtree, and lower than the nodes in the right subtree.

  \bigskip

  In the context of search trees, the (contents of the) nodes are
  called \textbf{keys}, and they need to be \textbf{totally ordered}
  (that is, the fact that any two keys must be comparable).

  \bigskip

  A consequence is that the inorder traversal of a binary search tree
  is the list of its keys sorted in non\hyp{}decreasing order.

\end{frame}

% ------------------------------------------------------------------------
%
\begin{frame}{Binary Search Trees}

  Here is an example:
  \begin{center}
    \includegraphics[bb=71 661 115 721]{bst_ex}
  \end{center}

  \smallskip

  The value for that tree is:
  \begin{tabbing}
    $t = \fun{Int}(17,$ \= $\fun{Int}(5,$ \= $\fun{Int} (3, \fun{Ext}(), \fun{Ext}()),$\\
                        \>                \> $\fun{Int}(11,$ \= $\fun{Ext}(),$\\
                        \>                \>                 \> $\fun{Int}(13, \fun{Ext}(), \fun{Ext}()))),$\\
                        \> $\fun{Int}(29, \fun{Ext}(), \fun{Ext}()))$
   \end{tabbing}

  \bigskip

  The inorder traversal is: $\fun{inorder} \; t \twoheadrightarrow [3;
    5; 11; 13; 17; 29].$

\end{frame}
